\scanpage{29}
\margin{Ici il semble \unsure{qu'on utilise} le fait que dans 1.1. on prend \(T\) \uline{projectif}}
\uuline{Théorème} 1.4. Soit \(X\) projective lisse connexe de dim \(n\) sur \(k\). \deleted{Alors pr les cycles de dim 0 et pr les cycles de codim 1, l'équiv.}

(i) Soit \(\nu\) un 0-cycle sur \(X\). Alors \(\nu\) est alb. équiv. à 0 ssi \(\nu\) est de degré 0, et \(S_X(\nu)=0\).

(ii) Soit \(D\) un diviseur sur \(X\). Alors \(D\) est alb-équi. à 0 ssi \(D\) est lin. équivalent à 0.

\margin{\uline{Remarque} 1.5. On ignore, \inserted{déjà} pr les 0-cycles, si l'équiv. d'Alb. est identique à l'équivalence linéaire. (Vrai pr tous les cycles???).}

\uuline{Démonstration}. (i) Si \(\nu\) est de degré 0 et \(S_X(\nu)=0\), on a \(\nu\) alb-équi. à 0 pr définition [prendre \(T=X\), \(Z=\text{diagonale de }X\), donc \(\nu=Z(\nu)\)]. L'assertion en sens inverse résulte de [\hspace{1em}, 1.2.]. Pour (ii), on sait déjà que lin. équi. à \(0\Rightarrow\) alb-équi. à 0, prouvons le sens inverse. On doit prouver que si
\mspage{28}
